% Generated by task tex-boundary from dossiers/boundary-adaptor.md, section E (E.1 to E.6), with
% the proposed restatements of its findings; blueprint excerpts copied from
% docs/roadmap/protocol-blueprint.md at b435631.
% Breakable inline code for this fragment's tables and prose (LuaLaTeX): \bndc{...} sets its
% argument verbatim in the monospace font, with break points after punctuation and between the
% words of a camel-case identifier. Defined once, whichever boundary fragment is read first.
\ifdefined\bndc\else
\directlua{
  bndbreak = function(s)
    local bs = string.char(92)
    local hash = string.char(35)
    s = s:gsub(hash .. hash, hash)
    local after = { ["."] = true, ["_"] = true, ["/"] = true, [":"] = true, [","] = true,
      ["-"] = true, [")"] = true, ["]"] = true, ["="] = true, ["|"] = true }
    local prev = ""
    for _, c in utf8.codes(s) do
      local ch = utf8.char(c)
      if ch:match("^[A-Z]$") and prev:match("^[a-z]$") then
        tex.sprint(bs .. "penalty400{}")
      end
      tex.sprint(-2, ch)
      if after[ch] or c > 127 then
        tex.sprint(bs .. "allowbreak{}")
      end
      prev = ch
    end
  end
}
\newcommand{\bndc}[1]{\texttt{\directlua{bndbreak("\luaescapestring{\detokenize{#1}}")}}}
\fi

\subsection{The adaptor: are its sketched statements well formed and provable?}\label{sec:gen-boundary-statements}

\begin{sloppypar}
Short forms in this section. \code{bp:} is the blueprint
\file{docs/roadmap/protocol-blueprint.md} at leanerVM \code{b435631}; every sketch below is
copied from it. Lean files are named as in \cref{sec:gen-boundary-relations}, at \code{b435631}.
Clean is cited at its old pin \code{93c9d1ef}, ArkLib at \code{dca90385}
(\file{Security/Basic.lean} and \file{Security/RoundByRound.lean} under
\file{ArkLib/OracleReduction/}), CompPoly at \code{3468b38c} (\file{MvPolyEquiv/Core.lean} and
\file{MvPolyEquiv/Eval.lean} under \file{CompPoly/Multivariate/}); Rust files are under
\file{crates/lean_vm/src/} at leanVM's pin \code{a386121f}, and \file{whir_config.rs} under
\file{crates/pcs/src/}. The base theorem of the proof system (T4) and end-to-end soundness (T7) are
named as in \cref{sec:gen-boundary-chain}.
\end{sloppypar}

The Layer 3 sketch (\bndc{bp:794-820}) and the Layer 2 sketch it needs (\bndc{bp:761-778}) are
read here as Lean to be written: for each sketched statement, whether it can be written, whether
it can be proved, and the restatement the review proposes.

\subsubsection{\texorpdfstring{\protect\code{witnessOf}}{witnessOf} and \texorpdfstring{\protect\code{satisfiedBy_witnessOf}}{satisfiedBy\_witnessOf}}\label{sec:gen-boundary-st-witnessof}

\begin{leancode}
def witnessOf (prog) (s) (q : Column (leanIsaInstance prog s).μ) :
    EnsembleWitness (leanIsaEnsemble prog)                            -- total and computable
theorem satisfiedBy_witnessOf (hs : s.Admissible prog)                 -- the caps: a hypothesis
    (h : M3Holds (leanIsaInstance prog s) input q) : SatisfiedBy prog input (witnessOf prog s q)
\end{leancode}
\src{docs/roadmap/protocol-blueprint.md:813-816 at b435631}

\paragraph{The definition.} \enquote{Total and computable} is achievable. \bndc{EnsembleWitness}
is in \bndc{Type 1} (probe \bndc{Shapes}, output lines 1--2); a computable definition may inhabit a
type in \bndc{Type 1} (the ensemble itself is a \bndc{def} that compiles), and the fields are
constructible: eight \bndc{Table}s with \bndc{component := (leanIsaEnsemble prog).tables[j]}, rows
built by one pass over the columns read through the layout,
\bndc{data := fun name n ↦ if name = "mem" ∧ n = 3 then … else #[]}, the public input off cells 0
and 1, and the three proof fields by \bndc{rfl}/\bndc{simp}. Total in \bndc{s}: for an inadmissible
\bndc{s} the tables are simply of the announced heights. Two things the sketch omits:
\bndc{witnessOf} has no \bndc{input} argument (\cref{tab:gen-boundary-entry}, the adaptor's row),
which is fine only if the lanes are read off \bndc{q} (then \bndc{public_input_eq} needs the first
two lines and \bndc{0 < logMem}); and the instance must exist first
(\cref{sec:gen-boundary-st-tables,sec:gen-boundary-st-poly}, \cref{f:boundary-flock-circular}).

\paragraph{The theorem.} It typechecks as a statement provided
\bndc{(leanIsaInstance prog s).Stmt} reduces to \bndc{PublicInput}. With a \bndc{def} instance the
elaborator unfolds it (probe \bndc{DefInstance}, the \bndc{example} at line 34 elaborates), but
instance search does not, so \bndc{Decidable (M3Holds (leanIsaInstance prog s) input q)}, which
every \bndc{#guard} of the Layer 3 tests needs, is found only for an \bndc{abbrev} instance (probe
\bndc{DefInstance}: the two guards on \bndc{instA} pass, the two on \bndc{inst} fail with
\enquote{Type mismatch … expected to have type Bool}; the spine's toy is an \bndc{abbrev} for the
same reason, \bndc{Spine/Toy.lean:94-95}). Provable as sketched: yes, with the additions of
\cref{sec:gen-boundary-conjuncts} (the boundary lemma, the count derivation, the Flock consequence
lemma).

\begin{change}{the adaptor's map and the instance's reducibility (\code{bp:803}, \code{bp:813})}
\asis \bndc{def witnessOf (prog) (s) (q : Column (leanIsaInstance prog s).μ)} and
\bndc{def leanIsaInstance (prog : Program) (s : Sizes) : M3Instance}.
\asproposed Either add \bndc{(input : PublicInput)} to \bndc{witnessOf} (the
\bndc{Refinement.map} has the statement) or say \enquote{the lanes are read off cells 0 and 1 of
\bndc{mem_0}, \bndc{mem_1}}. And: \enquote{\bndc{abbrev leanIsaInstance …}, reducible, so that
\bndc{Decidable (M3Holds (leanIsaInstance prog s) input q)} and the instances on \bndc{I.Stmt} are
found by instance search, as the toy is}.
\reason \Cref{f:boundary-witnessof-input-counts,f:boundary-reducible}.
\end{change}

\subsubsection{\texorpdfstring{\protect\code{stackOf}}{stackOf}, \texorpdfstring{\protect\code{m3Holds_stackOf}}{m3Holds\_stackOf}, \texorpdfstring{\protect\code{witnessOf_stackOf}}{witnessOf\_stackOf}}\label{sec:gen-boundary-st-stackof}

\begin{leancode}
def stackOf (gen : FlockWitnessGen) (w : EnsembleWitness (leanIsaEnsemble prog))
    (hs : Sizes.ofWitness w = some s) : Column (leanIsaInstance prog s).μ   -- gen: #3's
theorem m3Holds_stackOf (h : SatisfiedBy prog input w) (hs : Sizes.ofWitness w = some s) :
    M3Holds (leanIsaInstance prog s) input (stackOf gen w hs)
theorem witnessOf_stackOf (hs) : witnessOf prog s (stackOf gen w hs) = w  -- on the committed fields
\end{leancode}
\src{docs/roadmap/protocol-blueprint.md:811-812, 817-819 at b435631}

\bndc{stackOf} is well formed (a \bndc{def} taking a proof to fix the index of its result is
ordinary Lean). \bndc{m3Holds_stackOf} is provable with the additions of
\cref{sec:gen-boundary-completeness}.

\bndc{witnessOf_stackOf} is not a Lean statement. As an equality of \bndc{EnsembleWitness}es it is
false for almost every \bndc{w}: \bndc{w.data : String → (n : ℕ) → Array (Vector K n)} is a function
and \bndc{witnessOf} rebuilds it with the \bndc{"mem"} table only; \bndc{Table.width} is a free field
(\bndc{FlatComponent.lean:151-156}) and \bndc{witnessOf} rebuilds rows of exactly the component's
width; the memory block's \bndc{idx} cells and the bytecode block's entry cells are rebuilt from
\bndc{gpow} and \bndc{prog}, so equality needs \bndc{index_columns} and \bndc{bytecode_rows} of
\bndc{w}; and the public input is rebuilt from cells 0 and 1, so it needs \bndc{public_input_eq}
and the word conjuncts. \enquote{On the committed fields} has to be spelled out. Neither
composition of the base theorem uses \bndc{witnessOf_stackOf} (soundness uses
\bndc{satisfiedBy_witnessOf}, completeness \bndc{m3Holds_stackOf}); its role is acceptance tests 19
and 24 (\emph{The extractor reads the stack}, \emph{The extractor computes},
\bndc{bp:1320-1322, 1333-1340}), non-vacuity of \bndc{witnessOf}. The test
\bndc{#guard witnessOf (stackOf w) = w} (\bndc{bp:1339}) cannot be written, \bndc{EnsembleWitness}
having no \bndc{DecidableEq}.

\begin{change}{the round trip of the adaptor (\code{bp:819}, \code{bp:1339})}
\asis \bndc{theorem witnessOf_stackOf (hs) : witnessOf prog s (stackOf gen w hs) = w  -- on the committed fields}
and the test \bndc{#guard witnessOf (stackOf w) = w}.
\asproposed A statement that says what the sketch means and is decidable on a small witness:
\begin{leancode}
/-- The cells a component reads: the first `width` cells of a raw row, a missing cell `0`. -/
def rowCells (c : Component K) (row : Array K) : List K :=
  (List.range c.width).map fun j ↦ row[j]?.getD 0
theorem witnessOf_stackOf (h : SatisfiedBy prog input w) (hs : Sizes.ofWitness w = some s) :
    let w' := witnessOf prog s (stackOf gen w hs)
    (∀ j : Fin 8, (tableAt w' j).table.map (rowCells (tableAt w j).component) =
        (tableAt w j).table.map (rowCells (tableAt w j).component)) ∧
      memRows w'.data = memRows w.data ∧ w'.publicInput = w.publicInput
\end{leancode}
and, for \bndc{bp:1339}, \enquote{\bndc{#guard} of those three equalities on the one-row witness}.
Say in \bndc{bp:822-824} that neither composition of the base theorem uses it.
\reason \bndc{EnsembleWitness} has a function field, a free width and no \bndc{DecidableEq}
(\cref{f:boundary-witnessof-stackof}).
\end{change}

\subsubsection{The rate in \texorpdfstring{\protect\code{Sizes}}{Sizes}}\label{sec:gen-boundary-st-rate}

\begin{leancode}
structure Sizes where
  logMem : ℕ
  τ : Fin 6 → ℕ
  logInvRate : ℕ                                   -- WHIR rate exponent, announced with the sizes
def Sizes.ofWitness (w : EnsembleWitness leanIsaEnsemble) : Option Sizes   -- `none` unless every height is 2^τ
def Sizes.Admissible (prog : Program) (s : Sizes) : Prop      -- the caps; decidable
theorem admissible_iff_caps : s.Admissible prog ↔ (Caps w ∧ Sizes.ofWitness w = some s)
\end{leancode}
\src{docs/roadmap/protocol-blueprint.md:795-801 at b435631}

The Rust announces \bndc{log_mem}, six \bndc{τ_j} and \bndc{log_inv_rate} (\bndc{announce_public},
\bndc{cpu/mod.rs:118-124}; \bndc{read_public}, \bndc{cpu/mod.rs:144-149}); the Python reads the same
\bndc{2 + 6} scalars (\bndc{python-verifier/verifier.py:1372-1376}); the specification's Setup
names \bndc{κ_mem} and the six \bndc{τ_j} only
(\bndc{doc/leanvm/body/08-end-to-end-protocol.tex:55}), a disagreement between the specification
and both implementations that the blueprint resolves the Rust's way (its source category B,
\bndc{bp:199}). So the rate \emph{is} announced with the sizes. It does not belong in the instance
of the oracle protocol: \bndc{leanIsaInstance prog s} reads \bndc{s.logMem} and \bndc{s.τ} (heights,
\bndc{μ}, layout, lines) and nothing of the rate, which parameterizes WHIR alone (Layer 11);
\bndc{Sizes.ofWitness (w) : Option Sizes} (\bndc{bp:799}) cannot produce a rate from a witness; and
\bndc{Sizes.Admissible} would mix a check of the oracle protocol's family index with a check of
the commitment. The bus dossier reaches the same point from the caps
(\cref{f:bus-admissible}). Besides, \bndc{admissible_iff_caps} has a free \bndc{w} and, quantified,
is false in both directions (\bndc{Caps} lacks the \bndc{μ} and rate windows,
\bndc{Statement.lean:79-81}; another \bndc{w} has other sizes); \bndc{Sizes.ofWitness} lacks
\bndc{prog}; and \bndc{leanIsaInstance_fits} uses \bndc{layout.total} (\bndc{bp:808}), a field
\bndc{Layout} has not (the Layer 1 dossier, \file{code-layer1.md} G.2).

\begin{change}{the sizes, admissibility and the rate (\code{bp:795-808})}
\asis the sketch above, with \bndc{leanIsaInstance_fits : (leanIsaInstance prog s).layout.total ≤ 2 ^ (leanIsaInstance prog s).μ}
(\bndc{bp:808}).
\asproposed
\begin{leancode}
structure Sizes where (logMem : ℕ) (τ : Fin 6 → ℕ)                 -- the family index
def Sizes.ofWitness (w : EnsembleWitness (leanIsaEnsemble prog)) : Option Sizes
def leanIsaBlocks (prog) (s) : Blocks ; def leanIsaμ (prog) (s) : ℕ  -- witness.rs:67-101
theorem leanIsaBlocks_fits : (leanIsaBlocks prog s).total ≤ 2 ^ leanIsaμ prog s
def Sizes.Admissible (prog) (s) : Prop  -- 16 ≤ logMem ≤ 32, τ j ≤ 32, 3 ≤ τ 5, leanIsaμ prog s ≤ 28
def validRate (ρ : ℕ) : Prop := 1 ≤ ρ ∧ ρ ≤ 4                        -- whir_config.rs:48-55, Layer 12
theorem caps_of_admissible (hs : s.Admissible prog) (q) : Caps (witnessOf prog s q)
theorem sizes_of_satisfiedBy (h : SatisfiedBy prog input w) : ∃ s, Sizes.ofWitness w = some s
\end{leancode}
The announcement of Layer 12 is \bndc{Sizes × logInvRate}.
\reason The instance depends on the heights alone; the rate is Layer 12's; the two lemmas are what
the two directions use (\cref{f:boundary-sizes}).
\end{change}

\subsubsection{The instance's tables, columns and blocks against \texorpdfstring{\protect\code{M3Instance}}{M3Instance} as built}\label{sec:gen-boundary-st-tables}

\begin{leancode}
def leanIsaInstance (prog : Program) (s : Sizes) : M3Instance   -- `Ensemble.toM3` of the eight tables with
                                                               -- leanISA's separators and directions; the
                                                               -- layout of `witness.rs:67-101`, `leaf.rs:53-156`
theorem leanIsaInstance_degree (j) (C ∈ (leanIsaInstance prog s).constraints j) : C.totalDegree ≤ 2
theorem leanIsaInstance_flush_degree …  ≤ 2
theorem leanIsaInstance_fits : (leanIsaInstance prog s).layout.total ≤ 2 ^ (leanIsaInstance prog s).μ
\end{leancode}
\src{docs/roadmap/protocol-blueprint.md:803-808 at b435631}

and, in prose (\bndc{bp:836-839}): \enquote{The stack columns that belong to no opcode table
(\bndc{mem_0, mem_1, mem_2}, \bndc{cntfin_mem}, \bndc{cntfin_bc}, \bndc{q_flock};
\bndc{cpu/layout.rs:13-26}) are tables of the instance with no constraints and no flushes, so that
\bndc{Shape.ColumnId} and \bndc{Coord.committed} reach them}. Whether the instance the blueprint
means is \bndc{Ensemble.toM3} of the eight tables is \cref{f:boundary-not-ensemble-tom3}; read as
\cref{sec:gen-boundary-instance} reads it, it is expressible, with one consequence:
\begin{itemize}
\item the six shared columns as tables of the instance: yes, as at least three tables
  (\cref{sec:gen-boundary-instance}), each with \bndc{constraints := []}, \bndc{flushes := []},
  \bndc{counts := []}; \bndc{Shape.τ} is per table, so the memory columns' height \bndc{2^logMem}
  sits beside the six tables' heights;
\item the finalize counts outside the count tree: \bndc{counts jMem = []}, \bndc{counts jBc = []},
  as \bndc{layout.rs:412-414} (the count blocks are the tables' \bndc{count_columns()} only);
\item the boundary blocks: \bndc{Coord.committed ⟨jMem, i⟩ rfl} at \bndc{κ = logMem},
  \bndc{Coord.committed ⟨jBc, 0⟩ rfl} at \bndc{κ = prog.logSize}, the \bndc{known} and \bndc{const}
  coordinates of \cref{sec:gen-boundary-instance}; the \bndc{h : S.τ c.1 = κ} obligation is met by
  construction;
\item the eighteen limb columns as strided slots of \bndc{q_flock}: the \bndc{Layout} law admits it
  (\bndc{Spine/Instance.lean:73-81}, a reading law), but no reader exists (the Layer 1 dossier,
  \file{code-layer1.md} G.1) and the slot map is placed inside \bndc{FlockInterface (I)}
  (\cref{f:flock-ring-slot-map-circular}): \cref{f:boundary-flock-circular};
\item the three lines need \bndc{PublicLine.pos : 0 < τ}, that is \bndc{0 < s.logMem}; for a total
  \bndc{leanIsaInstance} the lines must be empty (or the instance junk) when \bndc{logMem = 0},
  which is harmless because \bndc{verify} rejects such \bndc{s}, but must be written;
\item \textbf{the consequence}: the generic table sumcheck (Layer 7) takes its rounds from
  \bndc{τ_max = max_j τ_j} over the instance's tables and sends \enquote{one value per column of
  every table} (\bndc{bp:987-989}, \bndc{08-end-to-end-protocol.tex:76-77}); with the shared columns
  as tables, a phase over the abstract \bndc{I} would run
  \bndc{max(τ_j, logMem, prog.logSize, τ_5 + 8)} rounds and send scalars for
  \bndc{mem_0, …, q_flock}, a transcript leanVM does not have (\bndc{constraints.rs:250}, the six
  tables; \cref{f:bus-table-phase-tables}). The instance needs a criterion the phase can read.
\end{itemize}

\begin{change}{column groups in the spine's conventions (\code{bp:316}) and in Layer 7 (\code{bp:987})}
\asis \enquote{tables of the instance with no constraints and no flushes} (\bndc{bp:836-839}),
with nothing saying how a phase over the abstract instance treats them.
\asproposed In the spine's conventions (\bndc{bp:316}, or a new row \emph{Column groups}):
\enquote{A table with no constraint, no flush and no count column is a column group: it takes no
part in the table sumcheck (no round, no term, no column value) and its columns receive claims
from the boundary blocks, the public lines and the Flock phase only.
\bndc{M3Instance.active j := (I.constraints j ≠ []) ∨ (I.flushes j ≠ []) ∨ (I.counts j ≠ [])} is the
criterion, decidable, and Layer 7 ranges over it.} In Layer 7 (\bndc{bp:987}): \bndc{τ_max} and
\bndc{Σ_j width_j} over the active tables.
\reason A phase over the abstract \bndc{I} that ranges over all of \bndc{I.ntab} has a different
transcript from leanVM's on the leanISA instance, which \bndc{verify_iff_compiled} would then fail
(\cref{f:boundary-column-groups}).
\end{change}

\subsubsection{Mathlib polynomials in Layer 2, CompPoly polynomials in the instance}\label{sec:gen-boundary-st-poly}

\begin{leancode}
def Expression.toMvPolynomial : Expression F → MvPolynomial ℕ F     -- var i ↦ X i
theorem eval_toMvPolynomial (env : Environment F) (e) :
    MvPolynomial.eval env.get e.toMvPolynomial = e.eval env
def Expression.degreeBound : Expression F → ℕ                      -- syntactic; add max, mul sum
theorem totalDegree_le_degreeBound (e) : e.toMvPolynomial.totalDegree ≤ e.degreeBound

structure M3Table (F) where
  width : ℕ
  constraints : List (MvPolynomial (Fin width) F)
  flushes : List (Direction × Vector (MvPolynomial (Fin width) F) 16)   -- separator first
  count : List (Fin width)                                             -- the count columns
def Component.toM3 (c : Component F) (sep : RawChannel F → F) (dir : RawChannel F → Direction) :
    M3Table F
theorem toM3_constraints_iff (t : Table F) (row ∈ t.table) :
    (∀ C ∈ (t.component.toM3 …).constraints, eval row C = 0) ↔ t.component.operations.ConstraintsHold (t.environment row)
theorem toM3_flushes_eq (t : Table F) : (multiset of evaluated flush tuples) = (multiset of `t.interactions` mapped to 16-tuples)
\end{leancode}
\src{docs/roadmap/protocol-blueprint.md:762-777 at b435631}

\begin{sloppypar}
Layer 2 produces \bndc{Expression.toMvPolynomial : Expression F → MvPolynomial ℕ F} and an
\bndc{M3Table F} whose constraints are \bndc{MvPolynomial (Fin width) F}; \bndc{M3Instance} holds
\bndc{CMvPolynomial (width j) K} (\bndc{Spine/Instance.lean:123-125}). No conversion is named.
CompPoly's \bndc{toCMvPolynomial : MvPolynomial (Fin n) R → CMvPolynomial n R} is
\bndc{noncomputable def} at the pin \bndc{3468b38c} (\bndc{MvPolyEquiv/Core.lean:41}; the same at the
new pin \bndc{572f9973}, read with \bndc{git show}), so an instance built through it is
noncomputable: \bndc{M3Holds (leanIsaInstance prog s)} could not be decided by evaluation (the
Layer 3 tests, \bndc{bp:866-867}), the phases' verifiers, which evaluate \bndc{I.constraints} at the
claimed column values, would not compile, and \bndc{verify} (\bndc{bp:1230-1231}, \enquote{total and
computable}) could not be defined from them. The degree bound travels the other way without
trouble: \bndc{totalDegree_equiv : p.totalDegree = (fromCMvPolynomial p).totalDegree}
(\bndc{MvPolyEquiv/Eval.lean:60-61}, \bndc{rfl}) and \bndc{eval_equiv} (\bndc{MvPolyEquiv/Eval.lean:53-57})
let Mathlib prove the bound of a computable polynomial, as the toy does
(\bndc{Spine/Toy.lean:68-74}).
\end{sloppypar}

What Layer 2 must produce is a \emph{direct} computable translation
\bndc{Expression K → CMvPolynomial n K} (structural recursion on \bndc{var}, \bndc{const},
\bndc{add}, \bndc{mul}, \bndc{Expression.lean:12-16}, with \bndc{CMvPolynomial.X}, \bndc{C},
\bndc{+}, \bndc{*}), and its bridge to Mathlib for the degree lemma. Probe \bndc{PolyBridge} writes it
in twelve lines, evaluates the \bndc{JUMP} residual \bndc{b + v_cond · w} on a row and checks it
against Clean's \bndc{Expression.eval}, checks \bndc{totalDegree = 2}, and checks that a variable
past the width reads \bndc{0} on both sides as \bndc{Environment.fromArray} reads a missing cell
(\bndc{Expression.lean:71-73}); it passed at the old pins, and as \bndc{PolyBridgeNew}, with
\bndc{K.ofBits} for the row numerals, at the new pins (\cref{sec:gen-boundary-probe-polybridge}).
The \bndc{Direction} of leanISA's \bndc{channelDir} (\bndc{Channels.lean:262-272}) and the spine's
\bndc{Side} (\bndc{Spine/Instance.lean:51-54}) are two enumerations of the same two values; the
translation is one match.

\begin{change}{the polynomial bridge of Layer 2 (\code{bp:762-767})}
\asis \bndc{def Expression.toMvPolynomial : Expression F → MvPolynomial ℕ F}, with \bndc{M3Table}'s
fields in \bndc{MvPolynomial (Fin width) F} and \bndc{Direction}.
\asproposed \enquote{\bndc{def Expression.toCMvPolynomial (n) : Expression K → CMvPolynomial n K}
(var \bndc{i < n} $↦$ \bndc{X i}, else \bndc{0}, as \bndc{Environment.fromArray} reads a missing cell);
\bndc{theorem eval_toCMvPolynomial (row) (e) : (e.toCMvPolynomial n).eval (fun i ↦ row[i]?.getD 0) = e.eval (Environment.fromArray row data)};
\bndc{theorem fromCMvPolynomial_toCMvPolynomial (e) : fromCMvPolynomial (e.toCMvPolynomial n) = rename … e.toMvPolynomial},
through which \bndc{degreeBound} bounds \bndc{totalDegree} (\bndc{totalDegree_equiv})}.
\bndc{M3Table}'s fields become \bndc{CMvPolynomial}; \bndc{Direction} becomes \bndc{Side} (or the
spine adopts leanISA's \bndc{Direction}). Clean \#466 (\bndc{Expression.toMvPolynomial}) stays the
proof-side bridge.
\reason A deviation forced by an upstream library, with the direct translation as its workaround,
retired if CompPoly makes \bndc{toCMvPolynomial} computable (\cref{f:boundary-poly-bridge}).
\end{change}

\subsubsection{The composition of the base theorem, arrow by arrow}\label{sec:gen-boundary-st-composition}

\begin{leancode}
   T4  =  verify_knowledgeSound  ∘  Refinement.map_option_valid satisfiedBy_witnessOf  ∘  constraintSoundness
\end{leancode}
\src{docs/roadmap/protocol-blueprint.md:407 at b435631}
\begin{leancode}
theorem baseVerifier_extractsExecution (fs bcs mca flock) (h : verify prog input proof = true) :
    except with probability niError, ∃ t, ValidExecution prog input t
theorem baseProver_complete (hfill : HasFillBlocks prog) (h : ValidExecution prog input t) :
    verify prog input (prove prog input (witness of t)) = true
\end{leancode}
\src{docs/roadmap/protocol-blueprint.md:1247-1250 at b435631}

\enquote{The first composes \bndc{verify_knowledgeSound}, the adaptor's pointwise transport
\bndc{Refinement.map_option_valid} with \bndc{satisfiedBy_witnessOf}, and \bndc{constraintSoundness}}
(\bndc{bp:1253-1254}); the tracker's section on the base theorem (K4) says
\bndc{knowledgeSound_of_refinement}, a name that exists nowhere. Arrow by arrow:
\begin{itemize}
\item \bndc{verify_knowledgeSound (fs bcs mca flock)}: in ArkLib's shape it bounds, for a fixed
  \bndc{stmtIn = input} and every prover, the probability of the event \enquote{the compiled
  verifier accepts and every \bndc{q} in the extracted slot has
  \bndc{(input, q) ∉ M3Rel (leanIsaInstance prog s)}} by \bndc{niError}
  (\bndc{Security/Basic.lean:299-323}, \bndc{knowledgeSoundnessWith}; the same event in the
  round-by-round form, \bndc{RoundByRound.lean:553}). The slot is an \bndc{Option (Column μ)},
  \bndc{Column μ : Type}.
\item \bndc{satisfiedBy_witnessOf}, made a \bndc{Refinement (M3Rel I) (SatRel prog)} with
  \bndc{SatRel prog := {p | SatisfiedBy prog p.1.1 p.2}} on witnesses in \bndc{Type 1}:
  \bndc{Refinement}'s two universes are independent (probe \bndc{Shapes}, output lines 10--13), and
  the construction typechecks (probe \bndc{Transport}, example 1).
  \bndc{Refinement.map_option_valid} has the polarity \enquote{every witness in the slot is valid
  $⇒$ every mapped witness is valid} (\bndc{ToArkLib/Refinement.lean:55-57}), while the game's bad
  event is \enquote{no witness in the slot is valid} (an empty slot is bad); the transport of the
  bad event is \bndc{∀ w' ∈ w?.map f, (x, w') ∉ S → ∀ w ∈ w?, (x, w) ∉ R}, which follows from
  \bndc{map_valid} in one line (probe \bndc{Transport}, \bndc{bad_of_bad}), and the spine dossier's
  probe \bndc{knowledge_transport} (\file{code-spine.md}) proves the probabilistic statement for
  ArkLib's game in the same way. \bndc{Extractor.Straightline.map} cannot serve: its \bndc{WitIn'}
  is in \bndc{Type} (probe \bndc{UniverseFail}: \enquote{\bndc{EnsembleWitness (leanIsaEnsemble prog)}
  has type \bndc{Type 1} … but is expected to have type \bndc{Type}}).
\item \bndc{constraintSoundness (hwf : WellFormedBytecode prog)}: pointwise.
\end{itemize}
The composition typechecks in principle as an inclusion of events (probe \bndc{Transport},
examples 3 and 4): if \bndc{¬ ∃ t, ValidExecution prog input t}, then every extracted \bndc{q} is
outside \bndc{M3Rel}, so the game's bad event is \enquote{the verifier accepts}, and
\bndc{Pr[accepts] ≤ niError}. That is the correct form of the first half of the base theorem:
\textbf{a soundness theorem for the language
\bndc{{(prog, input) | ∃ t, ValidExecution prog input t}}} (the conclusion is a closed
proposition, so \enquote{except with probability} attaches to acceptance, not to it), in the
random-oracle model, for each admissible \bndc{s} or with the family composition of
\cref{sec:gen-boundary-sizes-exists}, and under \bndc{WellFormedBytecode prog}. It is not, and cannot
be, an ArkLib \bndc{knowledgeSoundness} of a relation on \bndc{EnsembleWitness}
(\bndc{WitIn : Type}), nor a statement about the concrete \bndc{verify prog input proof = true} (a
closed Boolean), and it does not need the round-by-round implication if \bndc{FiatShamirSecurity}
consumes the round-by-round form directly.

\begin{change}{the base soundness theorem and the composition line (\code{bp:407}, \code{bp:1247-1248})}
\asis the two excerpts above.
\asproposed
\begin{leancode}
/-- Base soundness, in the random-oracle model: for a well-formed program and an input on
which no execution exists, every prover with at most `Q` oracle queries makes the
Fiat–Shamir-compiled verifier of `leanIsaInstance prog s` accept with probability at most
`niError Q` (the maximum over admissible `s`; finding 6). `verify` is that verifier with the
BLAKE2s chain in place of the oracle (`verify_iff_compiled`), which is a heuristic and not a
theorem. -/
theorem baseVerifier_sound (fs bcs mca flock) (hwf : WellFormedBytecode prog)
    (hno : ¬ ∃ t, ValidExecution prog input t) : ∀ prover, Pr[compiled verifier accepts] ≤ niError Q
/-- Pointwise, for the extractor's stack: what recursion consumes (T6). -/
theorem execution_of_extracted (hwf : WellFormedBytecode prog) (hs : s.Admissible prog)
    (h : M3Holds (leanIsaInstance prog s) input q) :
    ∃ t, AssignmentRepresents (witnessOf prog s q) t ∧ ValidExecution prog input t
\end{leancode}
(\enquote{finding 6} in the docstring is \cref{f:boundary-announced-sizes}.) The line at
\bndc{bp:407} becomes
\bndc{T4 = verify_knowledgeSound ∘ Refinement.knowledge_transport satisfiedBy_witnessOf ∘ constraintSoundness},
with \bndc{Refinement.knowledge_transport} (the event form of the spine dossier) named as what the
base theorem composes, and \bndc{map_option_valid} kept as the pointwise lemma of the honest
direction; the tracker's \bndc{knowledgeSound_of_refinement} the same.
\reason \Cref{f:boundary-base-soundness,f:boundary-refinement-polarity}.
\end{change}
